\documentclass[11pt,aspectratio=169]{beamer}
% ============================================================
% estilo_presentaciones.tex
% Preámbulo común de las presentaciones de Cálculo III.
%
% Uso: en cada .tex, justo después de \documentclass,
%
%     \documentclass[11pt,aspectratio=169]{beamer}
%     % ============================================================
% estilo_presentaciones.tex
% Preámbulo común de las presentaciones de Cálculo III.
%
% Uso: en cada .tex, justo después de \documentclass,
%
%     \documentclass[11pt,aspectratio=169]{beamer}
%     % ============================================================
% estilo_presentaciones.tex
% Preámbulo común de las presentaciones de Cálculo III.
%
% Uso: en cada .tex, justo después de \documentclass,
%
%     \documentclass[11pt,aspectratio=169]{beamer}
%     \input{estilo_presentaciones}
%     \setpie{Tema de la clase --- Cálculo III}
%     \setbloques{6}
%
% y en el cuerpo, \portada genera la diapositiva de título.
% ============================================================

% ============ PAQUETES ============
\usepackage[utf8]{inputenc}
\usepackage[T1]{fontenc}
\usepackage[spanish,es-tabla]{babel}
\usepackage{amsmath,amssymb,mathtools}
\usepackage{booktabs}
\usepackage{pgfplots}
\pgfplotsset{compat=1.17}
\usepackage{tikz}
\usetikzlibrary{arrows.meta,calc,decorations.pathmorphing,decorations.markings,positioning}
\usepackage{xcolor}

\DeclareMathOperator{\sen}{sen}

% OJO: con T1 el carácter «¿» literal se compone como «£» (en T1 el slot 191
% es la libra, no el signo de interrogación invertido). Escríbase siempre
% \textquestiondown{} y \textexclamdown{} en lugar de los caracteres literales.

% ============ COLORES ============
\definecolor{navy}{HTML}{1B2A4A}
\definecolor{tealx}{HTML}{1F7A6C}
\definecolor{brick}{HTML}{A34A3E}
\definecolor{gold}{HTML}{B8891C}
\definecolor{lightbg}{HTML}{F2F2EF}
\definecolor{gray60}{HTML}{8A8A85}

% ============ PARÁMETROS POR PRESENTACIÓN ============
% Texto del pie de página (izquierda).
\newcommand{\pietexto}{Cálculo III}
\newcommand{\setpie}[1]{\renewcommand{\pietexto}{#1}}
% Número total de bloques, usado por \bloqueframe.
\newcommand{\totalbloques}{6}
\newcommand{\setbloques}[1]{\renewcommand{\totalbloques}{#1}}

% ============ TEMA BASE ============
\usetheme{default}
\usefonttheme[onlymath]{serif}   % matemáticas en Computer Modern, texto en sans
\setbeamertemplate{navigation symbols}{}
\setbeamercolor{title}{fg=navy}
\setbeamercolor{frametitle}{bg=navy,fg=white}
\setbeamercolor{structure}{fg=tealx}
\setbeamercolor{block title}{bg=tealx,fg=white}
\setbeamercolor{block body}{bg=lightbg,fg=black}
\setbeamercolor{alerted text}{fg=brick}
\setbeamercolor{normal text}{fg=black}
\setbeamerfont{frametitle}{series=\bfseries}
\setbeamertemplate{frametitle}{
  \nointerlineskip
  \vskip-2pt
  \begin{beamercolorbox}[wd=\paperwidth,ht=1.15cm,dp=0.35cm,leftskip=0.5cm]{frametitle}
    \usebeamerfont{frametitle}\insertframetitle
  \end{beamercolorbox}
}
\setbeamertemplate{footline}{
  \leavevmode
  \hbox{\begin{beamercolorbox}[wd=\paperwidth,ht=0.5cm,dp=0.2cm,leftskip=0.5cm,rightskip=0.5cm]{structure}
    \footnotesize\color{gray60} \pietexto \hfill \insertframenumber/\inserttotalframenumber
  \end{beamercolorbox}}
}
\setbeamertemplate{itemize items}[circle]
\setbeamertemplate{blocks}[rounded][shadow=false]

% ============ DIAPOSITIVA DE TÍTULO ============
% Toma los datos de \title, \subtitle, \author, \institute y \date,
% que deben escribirse sin color: aquí se les aplica el de la portada.
\newcommand{\portada}{%
{
\setbeamertemplate{footline}{}
\begin{frame}[plain]
  \begin{tikzpicture}[remember picture,overlay]
    \fill[navy] (current page.south west) rectangle (current page.north east);
  \end{tikzpicture}
  \vspace{1.2cm}
  \centering
  {\color{white}\Huge\bfseries \inserttitle}\\[0.4em]
  {\color{tealx!60!white}\Large \insertsubtitle}\\[1.5em]
  {\color{white!85!black}\large \insertauthor}\\[0.2em]
  {\color{white!70!black}\normalsize \insertinstitute}\\[0.2em]
  {\color{white!60!black}\small \insertdate}
\end{frame}
}
}

% ============ CAJA DE NOTA DE ANIMACIÓN ============
\newcommand{\notaanimacion}[1]{%
  \begin{center}
    \fbox{\begin{minipage}{0.85\textwidth}\footnotesize\centering
    \textbf{\textcolor{gold}{$\blacktriangleright$ animación}}\quad #1
    \end{minipage}}
  \end{center}
}

% ============ CAJA DE IDEA CLAVE ============
\newcommand{\notaclave}[1]{%
  \begin{center}
    \fbox{\begin{minipage}{0.85\textwidth}\footnotesize\centering
    \textbf{\textcolor{tealx}{$\blacktriangleright$ clave}}\quad #1
    \end{minipage}}
  \end{center}
}

% ============ CAJA DE ADVERTENCIA ============
\newcommand{\alerta}[1]{%
  \begin{center}
    \fbox{\begin{minipage}{0.88\textwidth}\footnotesize\centering
    \textbf{\textcolor{brick}{$\blacksquare$ cuidado}}\quad #1
    \end{minipage}}
  \end{center}
}

% ============ FRAME DE DIVISOR DE BLOQUE ============
\newcommand{\bloqueframe}[3]{%
\begin{frame}[plain]
  \vfill
  \begin{center}
    {\color{gray60}\footnotesize BLOQUE #1 DE \totalbloques}\\[0.3em]
    {\color{navy}\Large\bfseries #2}\\[0.6em]
    {\color{tealx}\footnotesize #3}
  \end{center}
  \vfill
\end{frame}
}

% ============ TARJETA DE FIGURA (galerías) ============
\newcommand{\figcard}[2]{%
  \begin{minipage}[c]{0.48\textwidth}
    \centering
    #1\\[2pt]
    {\footnotesize #2}
  \end{minipage}
}

% ============ RASTREADOR DE PASOS ============
% Resalta el paso #3 dentro de una secuencia; cada presentación define
% sus propios rastreadores encima de \pasoitem.
\newcommand{\pasoitem}[3]{%
  \ifnum#1=#3
    \textbf{\textcolor{tealx}{#2}}%
  \else
    \textcolor{gray60}{#2}%
  \fi
}

%     \setpie{Tema de la clase --- Cálculo III}
%     \setbloques{6}
%
% y en el cuerpo, \portada genera la diapositiva de título.
% ============================================================

% ============ PAQUETES ============
\usepackage[utf8]{inputenc}
\usepackage[T1]{fontenc}
\usepackage[spanish,es-tabla]{babel}
\usepackage{amsmath,amssymb,mathtools}
\usepackage{booktabs}
\usepackage{pgfplots}
\pgfplotsset{compat=1.17}
\usepackage{tikz}
\usetikzlibrary{arrows.meta,calc,decorations.pathmorphing,decorations.markings,positioning}
\usepackage{xcolor}

\DeclareMathOperator{\sen}{sen}

% OJO: con T1 el carácter «¿» literal se compone como «£» (en T1 el slot 191
% es la libra, no el signo de interrogación invertido). Escríbase siempre
% \textquestiondown{} y \textexclamdown{} en lugar de los caracteres literales.

% ============ COLORES ============
\definecolor{navy}{HTML}{1B2A4A}
\definecolor{tealx}{HTML}{1F7A6C}
\definecolor{brick}{HTML}{A34A3E}
\definecolor{gold}{HTML}{B8891C}
\definecolor{lightbg}{HTML}{F2F2EF}
\definecolor{gray60}{HTML}{8A8A85}

% ============ PARÁMETROS POR PRESENTACIÓN ============
% Texto del pie de página (izquierda).
\newcommand{\pietexto}{Cálculo III}
\newcommand{\setpie}[1]{\renewcommand{\pietexto}{#1}}
% Número total de bloques, usado por \bloqueframe.
\newcommand{\totalbloques}{6}
\newcommand{\setbloques}[1]{\renewcommand{\totalbloques}{#1}}

% ============ TEMA BASE ============
\usetheme{default}
\usefonttheme[onlymath]{serif}   % matemáticas en Computer Modern, texto en sans
\setbeamertemplate{navigation symbols}{}
\setbeamercolor{title}{fg=navy}
\setbeamercolor{frametitle}{bg=navy,fg=white}
\setbeamercolor{structure}{fg=tealx}
\setbeamercolor{block title}{bg=tealx,fg=white}
\setbeamercolor{block body}{bg=lightbg,fg=black}
\setbeamercolor{alerted text}{fg=brick}
\setbeamercolor{normal text}{fg=black}
\setbeamerfont{frametitle}{series=\bfseries}
\setbeamertemplate{frametitle}{
  \nointerlineskip
  \vskip-2pt
  \begin{beamercolorbox}[wd=\paperwidth,ht=1.15cm,dp=0.35cm,leftskip=0.5cm]{frametitle}
    \usebeamerfont{frametitle}\insertframetitle
  \end{beamercolorbox}
}
\setbeamertemplate{footline}{
  \leavevmode
  \hbox{\begin{beamercolorbox}[wd=\paperwidth,ht=0.5cm,dp=0.2cm,leftskip=0.5cm,rightskip=0.5cm]{structure}
    \footnotesize\color{gray60} \pietexto \hfill \insertframenumber/\inserttotalframenumber
  \end{beamercolorbox}}
}
\setbeamertemplate{itemize items}[circle]
\setbeamertemplate{blocks}[rounded][shadow=false]

% ============ DIAPOSITIVA DE TÍTULO ============
% Toma los datos de \title, \subtitle, \author, \institute y \date,
% que deben escribirse sin color: aquí se les aplica el de la portada.
\newcommand{\portada}{%
{
\setbeamertemplate{footline}{}
\begin{frame}[plain]
  \begin{tikzpicture}[remember picture,overlay]
    \fill[navy] (current page.south west) rectangle (current page.north east);
  \end{tikzpicture}
  \vspace{1.2cm}
  \centering
  {\color{white}\Huge\bfseries \inserttitle}\\[0.4em]
  {\color{tealx!60!white}\Large \insertsubtitle}\\[1.5em]
  {\color{white!85!black}\large \insertauthor}\\[0.2em]
  {\color{white!70!black}\normalsize \insertinstitute}\\[0.2em]
  {\color{white!60!black}\small \insertdate}
\end{frame}
}
}

% ============ CAJA DE NOTA DE ANIMACIÓN ============
\newcommand{\notaanimacion}[1]{%
  \begin{center}
    \fbox{\begin{minipage}{0.85\textwidth}\footnotesize\centering
    \textbf{\textcolor{gold}{$\blacktriangleright$ animación}}\quad #1
    \end{minipage}}
  \end{center}
}

% ============ CAJA DE IDEA CLAVE ============
\newcommand{\notaclave}[1]{%
  \begin{center}
    \fbox{\begin{minipage}{0.85\textwidth}\footnotesize\centering
    \textbf{\textcolor{tealx}{$\blacktriangleright$ clave}}\quad #1
    \end{minipage}}
  \end{center}
}

% ============ CAJA DE ADVERTENCIA ============
\newcommand{\alerta}[1]{%
  \begin{center}
    \fbox{\begin{minipage}{0.88\textwidth}\footnotesize\centering
    \textbf{\textcolor{brick}{$\blacksquare$ cuidado}}\quad #1
    \end{minipage}}
  \end{center}
}

% ============ FRAME DE DIVISOR DE BLOQUE ============
\newcommand{\bloqueframe}[3]{%
\begin{frame}[plain]
  \vfill
  \begin{center}
    {\color{gray60}\footnotesize BLOQUE #1 DE \totalbloques}\\[0.3em]
    {\color{navy}\Large\bfseries #2}\\[0.6em]
    {\color{tealx}\footnotesize #3}
  \end{center}
  \vfill
\end{frame}
}

% ============ TARJETA DE FIGURA (galerías) ============
\newcommand{\figcard}[2]{%
  \begin{minipage}[c]{0.48\textwidth}
    \centering
    #1\\[2pt]
    {\footnotesize #2}
  \end{minipage}
}

% ============ RASTREADOR DE PASOS ============
% Resalta el paso #3 dentro de una secuencia; cada presentación define
% sus propios rastreadores encima de \pasoitem.
\newcommand{\pasoitem}[3]{%
  \ifnum#1=#3
    \textbf{\textcolor{tealx}{#2}}%
  \else
    \textcolor{gray60}{#2}%
  \fi
}

%     \setpie{Tema de la clase --- Cálculo III}
%     \setbloques{6}
%
% y en el cuerpo, \portada genera la diapositiva de título.
% ============================================================

% ============ PAQUETES ============
\usepackage[utf8]{inputenc}
\usepackage[T1]{fontenc}
\usepackage[spanish,es-tabla]{babel}
\usepackage{amsmath,amssymb,mathtools}
\usepackage{booktabs}
\usepackage{pgfplots}
\pgfplotsset{compat=1.17}
\usepackage{tikz}
\usetikzlibrary{arrows.meta,calc,decorations.pathmorphing,decorations.markings,positioning}
\usepackage{xcolor}

\DeclareMathOperator{\sen}{sen}

% OJO: con T1 el carácter «¿» literal se compone como «£» (en T1 el slot 191
% es la libra, no el signo de interrogación invertido). Escríbase siempre
% \textquestiondown{} y \textexclamdown{} en lugar de los caracteres literales.

% ============ COLORES ============
\definecolor{navy}{HTML}{1B2A4A}
\definecolor{tealx}{HTML}{1F7A6C}
\definecolor{brick}{HTML}{A34A3E}
\definecolor{gold}{HTML}{B8891C}
\definecolor{lightbg}{HTML}{F2F2EF}
\definecolor{gray60}{HTML}{8A8A85}

% ============ PARÁMETROS POR PRESENTACIÓN ============
% Texto del pie de página (izquierda).
\newcommand{\pietexto}{Cálculo III}
\newcommand{\setpie}[1]{\renewcommand{\pietexto}{#1}}
% Número total de bloques, usado por \bloqueframe.
\newcommand{\totalbloques}{6}
\newcommand{\setbloques}[1]{\renewcommand{\totalbloques}{#1}}

% ============ TEMA BASE ============
\usetheme{default}
\usefonttheme[onlymath]{serif}   % matemáticas en Computer Modern, texto en sans
\setbeamertemplate{navigation symbols}{}
\setbeamercolor{title}{fg=navy}
\setbeamercolor{frametitle}{bg=navy,fg=white}
\setbeamercolor{structure}{fg=tealx}
\setbeamercolor{block title}{bg=tealx,fg=white}
\setbeamercolor{block body}{bg=lightbg,fg=black}
\setbeamercolor{alerted text}{fg=brick}
\setbeamercolor{normal text}{fg=black}
\setbeamerfont{frametitle}{series=\bfseries}
\setbeamertemplate{frametitle}{
  \nointerlineskip
  \vskip-2pt
  \begin{beamercolorbox}[wd=\paperwidth,ht=1.15cm,dp=0.35cm,leftskip=0.5cm]{frametitle}
    \usebeamerfont{frametitle}\insertframetitle
  \end{beamercolorbox}
}
\setbeamertemplate{footline}{
  \leavevmode
  \hbox{\begin{beamercolorbox}[wd=\paperwidth,ht=0.5cm,dp=0.2cm,leftskip=0.5cm,rightskip=0.5cm]{structure}
    \footnotesize\color{gray60} \pietexto \hfill \insertframenumber/\inserttotalframenumber
  \end{beamercolorbox}}
}
\setbeamertemplate{itemize items}[circle]
\setbeamertemplate{blocks}[rounded][shadow=false]

% ============ DIAPOSITIVA DE TÍTULO ============
% Toma los datos de \title, \subtitle, \author, \institute y \date,
% que deben escribirse sin color: aquí se les aplica el de la portada.
\newcommand{\portada}{%
{
\setbeamertemplate{footline}{}
\begin{frame}[plain]
  \begin{tikzpicture}[remember picture,overlay]
    \fill[navy] (current page.south west) rectangle (current page.north east);
  \end{tikzpicture}
  \vspace{1.2cm}
  \centering
  {\color{white}\Huge\bfseries \inserttitle}\\[0.4em]
  {\color{tealx!60!white}\Large \insertsubtitle}\\[1.5em]
  {\color{white!85!black}\large \insertauthor}\\[0.2em]
  {\color{white!70!black}\normalsize \insertinstitute}\\[0.2em]
  {\color{white!60!black}\small \insertdate}
\end{frame}
}
}

% ============ CAJA DE NOTA DE ANIMACIÓN ============
\newcommand{\notaanimacion}[1]{%
  \begin{center}
    \fbox{\begin{minipage}{0.85\textwidth}\footnotesize\centering
    \textbf{\textcolor{gold}{$\blacktriangleright$ animación}}\quad #1
    \end{minipage}}
  \end{center}
}

% ============ CAJA DE IDEA CLAVE ============
\newcommand{\notaclave}[1]{%
  \begin{center}
    \fbox{\begin{minipage}{0.85\textwidth}\footnotesize\centering
    \textbf{\textcolor{tealx}{$\blacktriangleright$ clave}}\quad #1
    \end{minipage}}
  \end{center}
}

% ============ CAJA DE ADVERTENCIA ============
\newcommand{\alerta}[1]{%
  \begin{center}
    \fbox{\begin{minipage}{0.88\textwidth}\footnotesize\centering
    \textbf{\textcolor{brick}{$\blacksquare$ cuidado}}\quad #1
    \end{minipage}}
  \end{center}
}

% ============ FRAME DE DIVISOR DE BLOQUE ============
\newcommand{\bloqueframe}[3]{%
\begin{frame}[plain]
  \vfill
  \begin{center}
    {\color{gray60}\footnotesize BLOQUE #1 DE \totalbloques}\\[0.3em]
    {\color{navy}\Large\bfseries #2}\\[0.6em]
    {\color{tealx}\footnotesize #3}
  \end{center}
  \vfill
\end{frame}
}

% ============ TARJETA DE FIGURA (galerías) ============
\newcommand{\figcard}[2]{%
  \begin{minipage}[c]{0.48\textwidth}
    \centering
    #1\\[2pt]
    {\footnotesize #2}
  \end{minipage}
}

% ============ RASTREADOR DE PASOS ============
% Resalta el paso #3 dentro de una secuencia; cada presentación define
% sus propios rastreadores encima de \pasoitem.
\newcommand{\pasoitem}[3]{%
  \ifnum#1=#3
    \textbf{\textcolor{tealx}{#2}}%
  \else
    \textcolor{gray60}{#2}%
  \fi
}


\setpie{Longitud de arco --- Cálculo III}
\setbloques{6}

% ============ RASTREADOR DE PASOS (algoritmo de 4 pasos) ============
\newcommand{\pasotracker}[1]{%
  {\footnotesize\centering
  \pasoitem{1}{1. Parametrización}{#1}\ $\rightarrow$\
  \pasoitem{2}{2. Derivada}{#1}\ $\rightarrow$\
  \pasoitem{3}{3. Magnitud}{#1}\ $\rightarrow$\
  \pasoitem{4}{4. Integración}{#1}\par}
  \vspace{4pt}
}

% ============ RASTREADOR PARA REPARAMETRIZACIÓN ============
\newcommand{\stracker}[1]{%
  {\footnotesize\centering
  \pasoitem{1}{1. Derivada y magnitud}{#1}\ $\rightarrow$\
  \pasoitem{2}{2. Función $s(t)$}{#1}\ $\rightarrow$\
  \pasoitem{3}{3. Invertir $t(s)$}{#1}\ $\rightarrow$\
  \pasoitem{4}{4. Sustituir $\mathbf r(s)$}{#1}\par}
  \vspace{4pt}
}

\title{Longitud de arco}
\subtitle{Curvas parametrizadas en $\mathbb{R}^n$}
\author{Andrés Fabián Leal Archila}
\institute{Cálculo III (Multivariable) --- Cód. 20254}
\date{\today}

\begin{document}

\portada
% ---------- AGENDA ----------
\begin{frame}{Agenda de la clase}
  \begin{enumerate}
    \setlength{\itemsep}{10pt}
    \item \textbf{Teoría}: definición, fórmula, algoritmo \hfill {\small\color{gray60}(16 min)}
    \item \textbf{Ejemplo ancla}: la hélice circular \hfill {\small\color{gray60}(10 min)}
    \item \textbf{Un caso con más técnica}: parábola con sustitución \hfill {\small\color{gray60}(9 min)}
    \item \textbf{Otro caso técnico}: espiral logarítmica \hfill {\small\color{gray60}(10 min)}
    \item \textbf{Función longitud de arco}: reparametrización natural \hfill {\small\color{gray60}(13 min)}
    \item \textbf{Cierre}: cuándo falla la inversión, y curvas no suaves \hfill {\small\color{gray60}(9 min)}
  \end{enumerate}
  \vfill
  \begin{center}
    \small\color{gray60} Duración estimada total: $\sim$67 minutos
  \end{center}
\end{frame}

% ======================================================
% BLOQUE 1 --- TEORÍA
% ======================================================
\bloqueframe{1}{Teoría: definición y fórmula}{Motivación, teorema, algoritmo de 4 pasos}

\begin{frame}{Motivación: \textquestiondown{}qué es la longitud de una curva?}
  \begin{columns}
    \begin{column}{0.5\textwidth}
      \begin{itemize}
        \item Queremos medir la distancia recorrida \textbf{siguiendo la forma exacta} de la curva, no en línea recta.
        \pause
        \item Idea: aproximar la curva con una \textbf{poligonal} (segmentos rectos entre puntos de la curva).
        \pause
        \item Entre más segmentos usemos, mejor se ajusta la poligonal a la curva.
        \pause
        \item La longitud de arco es el \textbf{límite} de las longitudes poligonales cuando el número de segmentos $n\to\infty$.
      \end{itemize}
    \end{column}
    \begin{column}{0.45\textwidth}
      \centering
      \begin{tikzpicture}[scale=0.85]
        \draw[navy,thick,domain=0:3.4,samples=60] plot (\x,{1.6*sin(60*\x)+1.8});
        \foreach \x in {0,0.34,0.68,1.02,1.36,1.7,2.04,2.38,2.72,3.06,3.4}{
          \pgfmathsetmacro{\y}{1.6*sin(60*\x)+1.8}
          \fill[brick] (\x,\y) circle (1pt);
        }
        \draw[tealx,thick]
          (0,{1.6*sin(0)+1.8}) -- (0.34,{1.6*sin(60*0.34)+1.8})
          -- (0.68,{1.6*sin(60*0.68)+1.8}) -- (1.02,{1.6*sin(60*1.02)+1.8})
          -- (1.36,{1.6*sin(60*1.36)+1.8}) -- (1.7,{1.6*sin(60*1.7)+1.8})
          -- (2.04,{1.6*sin(60*2.04)+1.8}) -- (2.38,{1.6*sin(60*2.38)+1.8})
          -- (2.72,{1.6*sin(60*2.72)+1.8}) -- (3.06,{1.6*sin(60*3.06)+1.8})
          -- (3.4,{1.6*sin(60*3.4)+1.8});
      \end{tikzpicture}
      {\footnotesize\color{gray60} poligonal (teal) aproximando la curva (azul)}
    \end{column}
  \end{columns}
  \vspace{2pt}
  \notaanimacion{aproximación poligonal --- ver cómo converge al aumentar $n$}
\end{frame}

\begin{frame}{Definición formal}
  Sea $\mathbf r(t)=\langle x(t),y(t),z(t)\rangle$ con derivada continua en $[a,b]$.
  \pause
  \begin{block}{Teorema (Longitud de arco en $\mathbb R^3$)}
    \[
      L = \int_a^b \|\mathbf r'(t)\|\,dt
        = \int_a^b \sqrt{[x'(t)]^2+[y'(t)]^2+[z'(t)]^2}\,dt.
    \]
  \end{block}
  \pause
  \begin{block}{En $\mathbb R^2$}
    \[
      L = \int_a^b \sqrt{[x'(t)]^2+[y'(t)]^2}\,dt.
    \]
  \end{block}
  \pause
  \begin{center}
  \small\color{gray60} $\|\mathbf r'(t)\|$ es la \textbf{rapidez}: la fórmula dice que la distancia total es ``rapidez $\times$ tiempo'', integrada sobre todo el recorrido.
  \end{center}
\end{frame}

\begin{frame}{Conexión con Cálculo II}
  \begin{block}{Recordatorio}
    Si la curva es la gráfica de $y=f(x)$ (caso particular: $x=t$, $y=f(t)$), la fórmula se reduce a la que ya conoces:
    \[
      L=\int_a^b \sqrt{1+[f'(x)]^2}\,dx.
    \]
  \end{block}
  \vspace{6pt}
  \begin{center}
  \normalsize No es una fórmula nueva --- es la \textbf{misma idea}, extendida a curvas en $\mathbb R^3$ (o parametrizadas de cualquier forma, no solo $y=f(x)$).
  \end{center}
\end{frame}

\begin{frame}{Algoritmo: 4 pasos}
  \begin{center}
  \begin{tikzpicture}[node distance=0.4cm]
    \node[draw=tealx,thick,fill=tealx!10,rounded corners,minimum width=2.6cm,minimum height=1.2cm,align=center] (p1) {\bfseries 1\\Parametrización};
    \node[draw=tealx,thick,fill=tealx!10,rounded corners,minimum width=2.6cm,minimum height=1.2cm,align=center,right=of p1] (p2) {\bfseries 2\\Derivada};
    \node[draw=tealx,thick,fill=tealx!10,rounded corners,minimum width=2.6cm,minimum height=1.2cm,align=center,right=of p2] (p3) {\bfseries 3\\Magnitud};
    \node[draw=tealx,thick,fill=tealx!10,rounded corners,minimum width=2.6cm,minimum height=1.2cm,align=center,right=of p3] (p4) {\bfseries 4\\Integración};
    \draw[-Stealth,thick,gray60] (p1) -- (p2);
    \draw[-Stealth,thick,gray60] (p2) -- (p3);
    \draw[-Stealth,thick,gray60] (p3) -- (p4);
  \end{tikzpicture}
  \end{center}
  \vspace{10pt}
  \begin{enumerate}
    \item Identificar $\mathbf r(t)$ y el intervalo $[a,b]$.
    \item Calcular $\mathbf r'(t)$.
    \item Calcular $\|\mathbf r'(t)\|$.
    \item Evaluar $\displaystyle\int_a^b \|\mathbf r'(t)\|\,dt$.
  \end{enumerate}
  \vspace{6pt}
  \begin{center}\small\color{gray60} Este es el mapa que seguimos en cada ejemplo.\end{center}
\end{frame}

% ======================================================
% BLOQUE 2 --- HÉLICE (EJEMPLO ANCLA)
% ======================================================
\bloqueframe{2}{Ejemplo ancla: la hélice circular}{Un caso limpio, sin sustituciones}

\begin{frame}{Hélice circular --- planteamiento}
  \pasotracker{1}
  \begin{block}{Problema}
    Calcular la longitud de un arco de la hélice
    \[
      \mathbf r(t) = \langle \cos t,\ \sen t,\ t\rangle, \qquad 0\le t\le 2\pi.
    \]
  \end{block}
  \vspace{6pt}
  \begin{center}
  \small\color{gray60} Una vuelta completa alrededor del eje $z$, mientras sube una altura $2\pi$.
  \end{center}
\end{frame}

\begin{frame}{Hélice circular --- derivada}
  \pasotracker{2}
  Derivando componente a componente:
  \[
    \mathbf r'(t) = \langle -\sen t,\ \cos t,\ 1\rangle.
  \]
\end{frame}

\begin{frame}{Hélice circular --- magnitud}
  \pasotracker{3}
  \[
    \|\mathbf r'(t)\|
    = \sqrt{(-\sen t)^2+(\cos t)^2+1^2}
    = \sqrt{\sen^2 t+\cos^2 t+1}
    = \sqrt{1+1}=\sqrt 2.
  \]
  \pause
  \begin{center}
  \normalsize\color{tealx} El módulo es \textbf{constante}: la hélice se recorre a rapidez constante.
  \end{center}
\end{frame}

\begin{frame}{Hélice circular --- integración}
  \pasotracker{4}
  \[
    L = \int_0^{2\pi}\|\mathbf r'(t)\|\,dt = \int_0^{2\pi}\sqrt2\,dt = \sqrt2\,(2\pi-0) = 2\pi\sqrt2.
  \]
  \pause
  \[
    \boxed{L=2\pi\sqrt2}
  \]
\end{frame}

\begin{frame}{Hélice circular --- interpretación geométrica}
  \begin{itemize}
    \item La hélice da una vuelta completa mientras sube una altura $2\pi$.
    \item La longitud $2\pi\sqrt2$ es \textbf{mayor} que la proyección circular ($2\pi$) y mayor que la altura ($2\pi$).
    \pause
    \item Si se \textbf{desenrollara} la hélice sobre un plano, sería la hipotenusa de un triángulo rectángulo de catetos $2\pi$ y $2\pi$:
    \[
      \sqrt{(2\pi)^2+(2\pi)^2}=2\pi\sqrt2. \quad\checkmark
    \]
  \end{itemize}
  \vspace{8pt}
  \notaanimacion{hélice 3D + desenrollado al triángulo $2\pi\times2\pi$}
\end{frame}

% ======================================================
% BLOQUE 3 --- PARÁBOLA (SUSTITUCIÓN)
% ======================================================
\bloqueframe{3}{Un caso con más técnica}{Parábola $\langle t, t^2\rangle$: sustitución + fórmula de tabla}

\begin{frame}{Parábola --- planteamiento}
  \pasotracker{1}
  \begin{block}{Problema}
    Calcular la longitud de arco de $\mathbf r(t)=\langle t,\ t^2\rangle$, $0\le t\le 1$.
  \end{block}
  \vspace{6pt}
  \begin{center}\small\color{gray60} Esta curva va a volver a aparecer al final de la clase --- guárdala en mente.\end{center}
\end{frame}

\begin{frame}{Parábola --- derivada y módulo}
  \pasotracker{2}
  \[
    \mathbf r'(t)=\langle 1,\ 2t\rangle
  \]
  \pause
  \pasotracker{3}
  \[
    \|\mathbf r'(t)\| = \sqrt{1+4t^2}.
  \]
\end{frame}

\begin{frame}{Parábola --- integración}
  \pasotracker{4}
  \[
    L=\int_0^1\sqrt{1+4t^2}\,dt.
  \]
  \pause
  Sustitución $u=2t$, $du=2\,dt$:
  \[
    L = \frac12\int_0^2\sqrt{1+u^2}\,du.
  \]
  \pause
  Fórmula de tabla $\displaystyle\int\sqrt{1+u^2}\,du=\frac{u\sqrt{1+u^2}}{2}+\frac12\ln\!\left(u+\sqrt{1+u^2}\right)$:
  \[
    L=\frac12\left[\frac{u\sqrt{1+u^2}}{2}+\frac12\ln\!\left(u+\sqrt{1+u^2}\right)\right]_0^2.
  \]
\end{frame}

\begin{frame}{Parábola --- resultado}
  Evaluando en $u=2$ y $u=0$:
  \[
    L=\frac12\left[\sqrt5+\frac12\ln(2+\sqrt5)\right].
  \]
  \pause
  \[
    \boxed{L=\frac{\sqrt5}{2}+\frac14\ln\!\left(2+\sqrt5\right)\approx 1.479.}
  \]
\end{frame}

% ======================================================
% BLOQUE 4 --- ESPIRAL LOGARÍTMICA
% ======================================================
\bloqueframe{4}{Otro caso técnico}{Espiral $e^t\langle\cos t,\sen t\rangle$: álgebra limpia}

\begin{frame}{Espiral logarítmica --- planteamiento}
  \pasotracker{1}
  \begin{block}{Problema}
    Calcular la longitud de arco de $\mathbf r(t)=\langle e^t\cos t,\ e^t\sen t\rangle$, $0\le t\le\pi$.
  \end{block}
\end{frame}

\begin{frame}{Espiral logarítmica --- derivada}
  \pasotracker{2}
  Regla del producto en cada componente:
  \[
    \mathbf r'(t)=\langle e^t(\cos t-\sen t),\ e^t(\sen t+\cos t)\rangle.
  \]
\end{frame}

\begin{frame}{Espiral logarítmica --- módulo}
  \pasotracker{3}
  \begin{align*}
    \|\mathbf r'(t)\|^2
    &= e^{2t}(\cos t-\sen t)^2 + e^{2t}(\sen t+\cos t)^2\\
    \pause
    &= e^{2t}\bigl[(\cos^2t-2\sen t\cos t+\sen^2t)+(\sen^2t+2\sen t\cos t+\cos^2t)\bigr]\\
    \pause
    &= e^{2t}\cdot 2.
  \end{align*}
  \pause
  \begin{center}\normalsize\color{tealx} Los términos cruzados se cancelan.\end{center}
  \[
    \|\mathbf r'(t)\|=\sqrt2\,e^t.
  \]
\end{frame}

\begin{frame}{Espiral logarítmica --- integración}
  \pasotracker{4}
  \[
    L=\int_0^\pi \sqrt2\,e^t\,dt = \sqrt2\left[e^t\right]_0^\pi = \sqrt2\,(e^\pi-1).
  \]
  \pause
  \[
    \boxed{L=\sqrt2\,(e^\pi-1).}
  \]
\end{frame}

% ======================================================
% BLOQUE 5 --- FUNCIÓN LONGITUD DE ARCO
% ======================================================
\bloqueframe{5}{Función longitud de arco}{Reparametrización natural de una curva}

\begin{frame}{\textquestiondown{}Por qué reparametrizar por longitud de arco?}
  \begin{itemize}
    \item Hasta ahora, $t$ es solo un parámetro --- \textbf{no} necesariamente representa distancia recorrida.
    \pause
    \item Sería útil tener un parámetro que \textbf{sí} coincida con la distancia recorrida sobre la curva.
    \pause
    \item Eso es la \textbf{parametrización por longitud de arco}: útil en geometría diferencial (curvatura, torsión).
  \end{itemize}
  \pause
  \begin{block}{Definición}
    Sea $\mathbf r(t)$ suave ($\mathbf r'$ continua, $\mathbf r'(t)\neq\mathbf 0$) para $t\ge a$. La \textbf{función longitud de arco} desde $t=a$ es
    \[
      s(t)=\int_a^t \|\mathbf r'(u)\|\,du.
    \]
  \end{block}
\end{frame}

\begin{frame}{Invertibilidad}
  Por el Teorema Fundamental del Cálculo,
  \[
    \frac{ds}{dt}=\|\mathbf r'(t)\| > 0.
  \]
  \pause
  \begin{center}\normalsize\color{tealx} Como $\mathbf r$ es suave ($\mathbf r'(t)\neq\mathbf 0$), esta derivada es \textbf{siempre positiva}.\end{center}
  \pause
  \vspace{6pt}
  Por lo tanto $s(t)$ es \textbf{estrictamente creciente} $\Rightarrow$ \textbf{invertible}.
  \vspace{6pt}
  \pause
  \begin{center}\small\color{gray60} La hipótesis $\mathbf r'(t)\neq\mathbf 0$ es esencial --- volvemos a esto al final de la clase.\end{center}
\end{frame}

\begin{frame}{Reparametrización}
  Como $s(t)$ es invertible, existe $t(s)$, y podemos escribir
  \[
    \mathbf r(s) = \mathbf r(t(s)).
  \]
  \pause
  \begin{block}{Propiedad fundamental}
    \[
      \|\mathbf r'(s)\| = 1.
    \]
  \end{block}
  \pause
  \begin{center}
  \normalsize La rapidez es constante e igual a $1$: el parámetro $s$ \textbf{coincide} con la distancia recorrida.
  \end{center}
  \vspace{6pt}
  \begin{center}\small\color{gray60} Es la parametrización \emph{natural} de la curva --- elimina la dependencia de qué tan rápido se recorre.\end{center}
\end{frame}

\begin{frame}{Ejemplo: circunferencia --- planteamiento}
  \stracker{1}
  \begin{block}{Problema}
    Sea $\mathbf r(t)=\langle R\cos t,\ R\sen t\rangle$, $t\in[0,2\pi]$. Encontrar $s(t)$ y reparametrizar en términos de $s$.
  \end{block}
  \pause
  \[
    \mathbf r'(t)=\langle -R\sen t,\ R\cos t\rangle,
    \qquad
    \|\mathbf r'(t)\|=\sqrt{R^2\sen^2t+R^2\cos^2t}=R.
  \]
\end{frame}

\begin{frame}{Ejemplo: circunferencia --- función $s(t)$}
  \stracker{2}
  Tomando $a=0$:
  \[
    s(t)=\int_0^t R\,du = Rt.
  \]
  \pause
  \stracker{3}
  Despejando $t$:
  \[
    t=\frac{s}{R}.
  \]
\end{frame}

\begin{frame}{Ejemplo: circunferencia --- reparametrización y verificación}
  \stracker{4}
  Sustituyendo en $\mathbf r(t)$:
  \[
    \mathbf r(s)=\left\langle R\cos\frac{s}{R},\ R\sen\frac{s}{R}\right\rangle, \qquad 0\le s\le 2\pi R.
  \]
  \pause
  Verificación de rapidez unitaria:
  \[
    \mathbf r'(s)=\left\langle -\sen\frac{s}{R},\ \cos\frac{s}{R}\right\rangle,
    \qquad \|\mathbf r'(s)\|=1. \quad\checkmark
  \]
  \pause
  \[
    \boxed{s(t)=Rt,\qquad \mathbf r(s)=\left\langle R\cos\frac{s}{R},\ R\sen\frac{s}{R}\right\rangle}
  \]
\end{frame}

% ======================================================
% BLOQUE 6 --- CIERRE
% ======================================================
\bloqueframe{6}{Cierre}{Cuando no se puede invertir, y curvas no suaves}

\begin{frame}{La parábola, revisitada}
  Para $\mathbf r(t)=\langle t,\ t^2\rangle$, la función longitud de arco es
  \[
    s(t)=\int_0^t\sqrt{1+4u^2}\,du = \frac{t}{2}\sqrt{1+4t^2}+\frac14\ln\!\left(2t+\sqrt{1+4t^2}\right).
  \]
  \pause
  \begin{center}\normalsize\color{tealx} $s(t)$ \textbf{sí} tiene forma cerrada --- es el mismo cálculo del Bloque 3.\end{center}
  \pause
  \vspace{6pt}
  Pero \textbf{despejar $t$ en función de $s$} (invertir $s(t)$) no se puede hacer algebraicamente: mezcla un término lineal con un logaritmo.
  \pause
  \begin{center}
  \normalsize\color{brick}\textbf{$\mathbf r(s)$ queda solo definida implícitamente.}
  \end{center}
  \vspace{4pt}
  \begin{center}\small\color{gray60} No toda curva suave admite reparametrización explícita por longitud de arco.\end{center}
\end{frame}

\begin{frame}{El astroide: una curva no suave}
  \begin{block}{Curva}
    \[
      \mathbf r(t)=\langle \cos^3 t,\ \sen^3 t\rangle.
    \]
  \end{block}
  \pause
  \[
    \mathbf r'(t) = \langle -3\cos^2t\,\sen t,\ 3\sen^2t\,\cos t\rangle
    = 3\sen t\cos t\,\langle -\cos t,\ \sen t\rangle.
  \]
  \pause
  \begin{center}\normalsize\color{brick}\textbf{$\mathbf r'(t)=\mathbf 0$ en los múltiplos de $\pi/2$}\end{center}
  \vspace{6pt}
  \pause
  \begin{center}\small\color{gray60} Ahí la curva tiene \textbf{cúspides}: no es suave, y la función longitud de arco (tal como la definimos) no aplica en esos puntos.\end{center}
\end{frame}

\begin{frame}{El astroide: longitud de un arco}
  Para $0\le t\le \pi/2$ (donde sí es suave, salvo en los extremos):
  \[
    \|\mathbf r'(t)\| = 3|\sen t\cos t| = \tfrac32\sen(2t).
  \]
  \pause
  \[
    L=\int_0^{\pi/2} \tfrac32\sen(2t)\,dt = \left[-\tfrac34\cos(2t)\right]_0^{\pi/2}.
  \]
  \pause
  \[
    L = -\tfrac34\cos(\pi) - \left(-\tfrac34\cos 0\right) = \tfrac34+\tfrac34.
  \]
  \pause
  \[
    \boxed{L=\tfrac32.}
  \]
\end{frame}

\begin{frame}{Resumen}
  \begin{itemize}
    \item Longitud de arco: $\displaystyle L=\int_a^b\|\mathbf r'(t)\|\,dt$ --- misma idea de Cálculo II, extendida a $\mathbb R^n$.
    \item El cálculo siempre sigue el mismo algoritmo de 4 pasos.
    \item La función $s(t)=\int_a^t\|\mathbf r'(u)\|\,du$ permite reparametrizar por longitud de arco, con $\|\mathbf r'(s)\|=1$.
    \item Esto requiere $s(t)$ invertible (curva suave, $\mathbf r'\neq\mathbf 0$) --- y aun siendo invertible en teoría, no siempre se puede despejar $t(s)$ en forma cerrada.
  \end{itemize}
  \vspace{10pt}
  \pause
  \begin{block}{Próxima clase}
    Con la parametrización por longitud de arco lista, podemos definir \textbf{curvatura} y \textbf{torsión}: cuánto se dobla y cuánto se tuerce una curva en el espacio.
  \end{block}
\end{frame}

\end{document}
